\documentclass[10pt,a4paper]{amsart}
\usepackage[margin=19mm]{geometry}
\usepackage{amsmath,amssymb,mathtools}
\usepackage[scr=rsfs,cal=euler]{mathalfa}
\usepackage{xfrac}
\usepackage[unicode,hidelinks,bookmarks=false]{hyperref}
\newcommand{\ie}{i.e.\ }
\newcommand{\cf}{cf.\ }
\newcommand{\resp}{respectively\ }
\newcommand{\Subsection}{\subsection}
\providecommand{\nobreakdash}{\nobreak\hbox{-}}
\setlength{\emergencystretch}{2em}
\multlinegap0pt
\usepackage{bm}
\usepackage{bbm}
\usepackage{dsfont}
\usepackage{xspace}
\usepackage{cancel}
\usepackage{mathtools}
\usepackage{amsthm}
\makeatletter
\@ifundefined{theorem}{\newtheorem{theorem}{Theorem}}{}
\@ifundefined{lemma}{\newtheorem{lemma}[theorem]{Lemma}}{}
\makeatother
\def\de{{\rm d}}
\title[Conditioned Brownian motion and local equivalence of path en]{Conditioned Brownian motion and local equivalence of path ensembles}
\author{Tobias Schmidt}
\date{Source: arXiv:2608.19744v1 (CC BY 4.0)}
\begin{document}
\maketitle

\noindent\emph{Source and licence.} Conditioned Brownian motion and local equivalence of path ensembles. Version arXiv:2608.19744v1 (CC BY 4.0), distributed under the Creative Commons Attribution 4.0 International licence, as recorded in the arXiv licence metadata for this exact version and shown on the abstract page https://arxiv.org/abs/2608.19744v1. Full rights notice: licenses/arxiv-2608.19744-CC-BY-4.0.txt.\par\smallskip

\noindent\textit{Editorial scope.} Counted unit: the Lemma whose source label is \texttt{lem:range-constraint-coupling} in the article (source0.tex lines 3435--3451, source label \texttt{lem:range-constraint-coupling}), delivered together with the article's own complete original proof of it (source0.tex lines 3453--3555, one \texttt{proof} environment of 103 lines). No mathematics is written, repaired or re-derived: statement and proof are the article's own text line for line. Required context retained uncounted: none. Every definition, display and label of those lines is retained unchanged. Omitted from this package and not redistributed: 1--3434, 3452--3452, 3556--3733. Nothing inside a retained excerpt is omitted or altered: every retained body is delivered exactly as the article's lines are written, so no omission receipt of any kind accompanies this package. Citations: the retained excerpts carry no citation command at all, so no bibliography entry of the article is transcribed and no reference is redistributed. Reference adaptation: none was needed and none was made; every reference inside the delivered unit resolves inside it, so no unresolved reference or citation is produced. Printed slips kept literal: no printed word, symbol, hypothesis or number of the counted statement or of its proof was altered. Layout and class: the article's own class is \texttt{amsart} with packages including xcolor, comment, hyperref, cleveref, geometry, newtxtext, newtxmath, parskip; the wrapper is the lane's pinned \texttt{amsart} editorial wrapper at 10pt on A4, which restates the theorem-like environments the retained text needs and adds the article's own macro definitions that the retained text needs (\texttt{\textbackslash de}), with extra wrapper packages (bm, bbm, dsfont, xspace, cancel, mathtools). The delivered numbering is the unit's own. Assets: no retained span contains a figure environment, a graphics-inclusion command, a PDF-page inclusion, a table or a TikZ picture; no image file is copied into this package, the package declares no asset dependency and no image file is redistributed with it. Licence: version arXiv:2608.19744v1 of this article is distributed under the Creative Commons Attribution 4.0 International licence, as recorded in the arXiv licence metadata for this exact version and shown on the abstract page https://arxiv.org/abs/2608.19744v1. A full rights notice is delivered at licenses/arxiv-2608.19744-CC-BY-4.0.txt. Characters: every retained excerpt is pure ASCII, so the delivered engine, XeLaTeX with its default Latin Modern text font, reports no missing character.
\begin{lemma}
\label{lem:range-constraint-coupling}
Assume that $Q\colon\mathbb R^d\to[0,\infty)$ is continuous,
$Q(x)\to\infty$ as $|x|\to\infty$, and is not constant.  Then
\[
    \lim_{\beta\downarrow0}\lambda'(\beta)=\infty,
    \qquad
    \lim_{\beta\uparrow\infty}\lambda'(\beta)=q_*.
\]
Consequently,
\[
    \beta\longmapsto\lambda'(\beta)
\]
is a continuous strictly decreasing bijection from $(0,\infty)$ onto
$(q_*,\infty)$.  In particular, for every $\theta>q_*$ there is a unique
$\beta_\theta>0$ such that $\lambda'(\beta_\theta)=\theta$.
\end{lemma}

\begin{proof}
Continuity and strict monotonicity of $\lambda'$ follow from the real
analyticity of $\lambda$ and the strict inequality $\lambda''(\beta)<0$
proved above.  Moreover,
\[
    \lambda'(\beta)-q_*
    =\int_{\mathbb R^d}\de x \, (Q(x)-q_*)\phi_\beta(x)^2>0.
\]
Indeed, $Q-q_*\geq0$, the ground state is strictly positive, and $Q$ is not
constant.  It remains to identify the two endpoint limits.
We first prove that
\begin{equation}
    \lambda(\beta)\longrightarrow0
    \qquad\text{as }\beta\downarrow0.
\label{eq:lambda-small-beta}
\end{equation}
Choose $\eta\in C_c^\infty(B(0,1))$ with $\|\eta\|_2=1$ and put
$f_L(x)=L^{-d/2}\eta(x/L)$.
Then $\|f_L\|_2=1$ and
\[
    \frac12\|\nabla f_L\|_2^2
    =\frac{1}{2L^2}\|\nabla\eta\|_2^2,
    \qquad
    \int_{\mathbb R^d} \de x \, Q(x)|f_L(x)|^2
    \leq \sup_{|x|\leq L}Q(x)<\infty.
\]
By Rayleigh-Ritz,
\[
    0\leq\lambda(\beta)
    \leq \frac{1}{2L^2}\|\nabla\eta\|_2^2
       +\beta\sup_{|x|\leq L}Q(x).
\]
Given $\varepsilon>0$, we first choose $L$ so that the kinetic term is at
most $\varepsilon/2$, and then choose $\beta$ so that the potential term is
at most $\varepsilon/2$.  This proves \eqref{eq:lambda-small-beta}.

Suppose now, contrary to the first claimed limit, that there are
$\beta_n\downarrow0$ and $M<\infty$ such that $\theta_{\beta_n}=\lambda'(\beta_n)\leq M$ for every $n$.
Write $\phi_n=\phi_{\beta_n}$.  Taking the inner product of the
ground-state equation with $\phi_n$ gives
\[
    \lambda(\beta_n)
    =\frac12\|\nabla\phi_n\|_2^2
      +\beta_n\theta_{\beta_n}.
\]
In view of \eqref{eq:lambda-small-beta} it must hold that
\begin{equation}
    \|\nabla\phi_n\|_2\longrightarrow0.
\label{eq:derivative-vanishes}
\end{equation}
On the other hand, the bound
\[
    \int_{\mathbb R^d}\de x \, Q(x) \phi_n(x)^2 \leq M
\]
makes the probability measures $\phi_n(x)^2\,\mathrm dx$ tight, because
$Q(x)\to\infty$.  In particular, one can choose $R<\infty$, independently
of $n$, such that
\begin{equation}
    \int_{B(0,R)} \de x \, \phi_n(x)^2\geq\frac12.
\label{eq:tight-mass}
\end{equation}
For each $m\geq R$, the sequence $(\phi_n)$ is bounded in $H^1(B(0,m))$.
By Rellich compactness and a diagonal subsequence, there is
$\phi\in L^2_{\mathrm{loc}}(\mathbb R^d)$ such that
$\phi_n\to\phi$ strongly in $L^2(B(0,m))$ for every $m$.  Equation
\eqref{eq:derivative-vanishes} implies $\nabla\phi=0$ in the distributional
sense.  Since $\mathbb R^d$ is connected, $\phi=c$ almost everywhere for
some constant $c$.  The mass bound \eqref{eq:tight-mass} gives
$|c|^2|B(0,R)|\geq1/2$, so $c\neq0$.  On the other hand, local $L^2$
convergence and $\|\phi_n\|_2=1$ give
$|c|^2|B(0,m)|\leq1$ for every $m$, a contradiction as $m\to\infty$.
Hence
$\lambda'(\beta)\to\infty$ as $\beta\downarrow0$.
We finally consider $\beta\to\infty$.  Fix $\varepsilon>0$ and choose
$x_*\in\mathbb R^d$ with $Q(x_*)=q_*$.  By continuity, there is a nonempty
open neighborhood $U_\varepsilon$ containing $x_*$ such that $Q(x)\leq q_*+\varepsilon$ for $x\in U_\varepsilon$.
Choose $f_\varepsilon\in C_c^\infty(U_\varepsilon)$ with
$\|f_\varepsilon\|_2=1$ and set $K_\varepsilon=\frac12\|\nabla f_\varepsilon\|_2^2$.
Rayleigh-Ritz gives
\begin{equation}
    \lambda(\beta)
    \leq K_\varepsilon+\beta(q_*+\varepsilon).
\label{eq:large-beta-upper}
\end{equation}
Since $Q\geq q_*$, we also have
\begin{equation}
    \lambda(\beta/2)\geq\frac\beta2q_*.
\label{eq:large-beta-lower}
\end{equation}
The function $\lambda$ is concave, being the infimum of affine functions of
$\beta$.  Therefore its derivative at the right endpoint is bounded above by
\[
    \lambda'(\beta)
    \leq \frac{\lambda(\beta)-\lambda(\beta/2)}{\beta/2}.
\]
Using \eqref{eq:large-beta-upper} and \eqref{eq:large-beta-lower}, we obtain
\[
    q_*<\lambda'(\beta)
    \leq q_*+2\varepsilon+\frac{2K_\varepsilon}{\beta}.
\]
Taking first $\beta\to\infty$ and then $\varepsilon\downarrow0$ proves $\lim_{\beta\to\infty}\lambda'(\beta)=q_*$.
The endpoint limits, continuity, and strict monotonicity complete the proof.
\end{proof}

\end{document}
